\documentclass{article}
\usepackage[UTF8]{ctex}
\usepackage{amsmath,amssymb,amsfonts}
\usepackage{enumitem}
\usepackage{geometry}
\geometry{a4paper, margin=1in}
% === 错题本排版优化（不参与题目标识） ===
\setlength{\parindent}{0pt}
\setlength{\parskip}{0.35em plus 0.10em minus 0.08em}
\setlength{\emergencystretch}{3em}
\setlist{itemsep=0.25em, topsep=0.40em, parsep=0pt}
\allowdisplaybreaks[2]
% === 错题本排版优化结束 ===
\begin{document}
\section*{首都师范大学2024年高等代数真题}
\section*{28．首都师范大学2024年高等代数试卷}

\subsection{题目1}
1．设 \(n\) 为正整数，计算 \(n\) 阶行列式

\[
\left|\begin{array}{cccc}
1+x_{1} & x_{1} & \cdots & x_{1} \\
x_{2} & 1+x_{2} & \cdots & x_{2} \\
\vdots & \vdots & & \vdots \\
x_{n} & x_{n} & \cdots & 1+x_{n}
\end{array}\right|
\]

\subsection{题目2}
2．判定下列多项式是否为有理数域上的不可约多项式，并说明理由。

\begin{itemize}
  \item[（1）] \(x^{8}+x^{4}+1\) ．
  \item[（2）] \(x^{5}+5 x+11\) ．
\end{itemize}

\subsection{题目3}
3．设某向量空间中的五个向量 \(\alpha_{1}, \alpha_{2}, \cdots, \alpha_{5}\) 满足 \(\alpha_{3}=\alpha_{1}-\alpha_{2}-\alpha_{5}, \alpha_{4}=\alpha_{1}-\alpha_{3}+\alpha_{5}\) ，且这组向量的秩 \(r\left(\alpha_{1}, \alpha_{2}, \cdots, \alpha_{5}\right)=3\) ．求方程组 \(x_{1} \alpha_{1}+x_{2} \alpha_{2}+\cdots+x_{5} \alpha_{5}=0\) 的一般解．

\subsection{题目4}
4．证明：不存在复数域上的方阵 \(A\) ，使得 \(A^{2}=\left(\begin{array}{lll}0 & 0 & 0 \\ 1 & 0 & 0 \\ 0 & 1 & 0\end{array}\right)\) ．

\subsection{题目5}
5．证明：奇数维欧氏空间的任一正交变换的全部实特征值之和必为奇数．

\subsection{题目6}
6．设 \(A=\left(\begin{array}{cc}1 & 2 \\ 2023 & 2022\end{array}\right)\) ，计算 \(A^{2024}\) ．

\subsection{题目7}
7．设 \(\mathbb{C}\) 为复数域，\(n\) 为大于1的整数，\(V=\left\{\sum_{0 \leq i<n} c_{i} x^{i} \mid c_{i} \in \mathbb{C}\right\}\) 为多项式环 \(\mathbb{C}[x]\) 的一个子集，定义映射 \(\mathcal{D}: V \rightarrow V, \sum_{0 \leq i<n} c_{i} x^{i} \mapsto \sum_{1 \leq i<n} i c_{i} x^{i-1}\) ．

\begin{itemize}
  \item[（1）] 证明：\(V\) 是 C 上的一个线性空间， \(\mathcal{D}\) 是 \(V\) 上的一个线性变换。
  \item[（2）] 求 \(\mathcal{D}\) 的极小多项式和特征多项式。
  \item[（3）] 求 \(\mathcal{D}\) 的所有不变子空间．
\end{itemize}

\subsection{题目8}
8．设 \(n\) 为正整数，\(M(n, \mathbb{R})\) 为实数域 \(\mathbb{R}\) 上所有 \(n\) 阶矩阵构成的线性空间，给定两个方阵 \(A, B \in M(n, \mathbb{R})\) ，定义映射 \(\varphi: M(n, \mathbb{R}) \rightarrow M(n, \mathbb{R}), X \mapsto A X B\) 。

\begin{itemize}
  \item[（1）] 证明：\(\varphi\) 为线性变换．
  \item[（2）] 设 \(a=\operatorname{tr}(A), b=\operatorname{tr}(B)\) 分别为矩阵 \(A, B\) 的迹（一个方阵的迹等于这个矩阵对角线上的元素之和），求 \(\operatorname{tr}(\varphi)\) ．
  \item[（3）] 设 \(|A|,|B|\) 分别为 \(A, B\) 的行列式，求 \(\varphi\) 的行列式。
\end{itemize}

\begin{itemize}

\subsection{题目9}
\end{itemize}

\subsection{题目9}
原题边界“9．”：利用正交变换把下面的实二次型化为标准型。
\[
7 x_{1}^{2}+7 x_{2}^{2}-5 x_{3}^{2}+32 x_{1} x_{2}-16 x_{1} x_{3}-16 x_{2} x_{3} .
\]
\end{document}
